\documentclass[11pt]{article}
\usepackage[margin=25mm]{geometry}
\usepackage{amsmath,amssymb,amsthm,booktabs,hyperref}
\newtheorem{theorem}{Theorem}\newtheorem{lemma}{Lemma}
\newcommand{\C}{\mathrm C}\newcommand{\A}{\mathrm A}\newcommand{\Sym}{\mathrm S}\newcommand{\PSL}{\mathrm{PSL}}
\title{Group certificates for twenty-two hyperbolic table entries}
\author{Finite companion to the fixed-source recovery}\date{}
\begin{document}\maketitle
\begin{abstract}
We identify actual automorphism groups of twenty-two marked quotients already
certified as regular rank-four hypertopes. Explicit homomorphisms, normal
kernels and complements establish the group descriptions in Table~3 of
Fernandes--Leemans--Weiss. The proof uses no new normal-subgroup search.
The two excluded orders and the ninth diagram remain separate obligations
of the complete source contract.
\end{abstract}
\section{The precise assertion}
Put $E_d=\C_2^d$, $L_q=\PSL_2(q)$, and let $X_{32}$ denote the actual
extraspecial group of order32 certified below. A colon means a split
extension, a cross a direct product, and a dot a nonsplit extension.
An expression with an unspecified action denotes the actual extension
supplied by its certificate; it is not a uniqueness assertion about every
group with that order or series.

We acquire the complete eight-type quotient packet and its distinct review
and root adoption. Its record $(\tau,j)$ is an ordered four-tuple of
permutations; $H_{\tau,j}$ is the literal group generated by that tuple.
The packet proves the exact Coxeter orders, all256 subset intersections and
actual flag transitivity for every selected record. Thus these properties
are not inferred here from the abstract group name. Its numbering is fixed;
a freshly enumerated kernel list need not have the same order.

\begin{theorem}
For each row of Table~\ref{tab}, the actual marked group $H_{\tau,j}$ has the
displayed description. Consequently all twenty-two positive named entries
of fixed-source Table~3 outside the diagram $(3,5,3,5)$ and the two excluded
orders are realized by the previously certified hypertopes.
\end{theorem}
The first four digits denote successive square edges; three digits denote
a string. The label $5;3,3$ denotes the star with three edges at vertex1.
The source's ``locally spherical'' diagram condition and the stronger
universal-residue condition remain distinct, as in the acquired core.
The source group notation need not uniquely fix an action. For these actual
twenty-two records, the stored local orders equal the universal orders in
the acquired core. Thus these witnesses also have universal rank-three
residues; this follows from their exact local data, not their group names.

\section{Certificate principles}
For finite permutation groups $N,K\leq D$, the equations
\[
 N\triangleleft D,\qquad N\cap K=1,\qquad \langle N,K\rangle=D
\]
give $D=N\rtimes K$. A specified bijective homomorphism $K\to M$ identifies
the complement with the named model $M$. An involution $t\notin D$ with
$D\triangleleft\langle D,t\rangle$ and index2 gives another split step.
If $z$ is a central involution outside the preceding subgroup $U$, then
$\langle U,z\rangle=U\times\langle z\rangle$.
These statements follow immediately from the unique product expression in
subgroups with trivial intersection; they require actual containment and
normality, which the checker tests explicitly.
\newpage
\begin{table}[t]\centering\small
\begin{tabular}{@{}lrrl@{}}\toprule
Diagram $\tau$&Record $j$&Order&Certified description\\\midrule
$(3,3,3,4)$&8&1920&$E_4:\A_5:\C_2$\\
&9&2688&$(E_3:L_7)\times\C_2$\\
&11&3840&$E_5:\A_5:\C_2$\\
&12&7680&$(\C_4\mathbin{.}(E_4:\A_5)):\C_2$\\
&13&7680&$E_6:\A_5:\C_2$\\
$(3,3,3,5)$&5&7200&$(\A_5\times\A_5):\C_2$\\
$(3,4,3,4)$&24&1152&$X_{32}:\Sym_3:\Sym_3$\\
&27&1440&$\Sym_6\times\C_2$\\
&36&4320&$\Sym_6\times\Sym_3$\\
&38&5376&$(E_3:L_7)\times\C_2^2$\\
&40&6048&$L_7\times\Sym_3^2$\\
&45&9216&$E_8:\Sym_3:\Sym_3$\\
$[3,5,3]$&5&8160&$L_{16}:\C_2$\\
$[4,3,5]$&14&3840&$E_5:\A_5:\C_2$\\
&12&3840&$(E_4:\A_5:\C_2)\times\C_2$\\
&15&7680&$(E_5:\A_5:\C_2)\times\C_2$\\
&16&7680&$E_5:\A_5:\C_2:\C_2$\\
$[5,3,5]$&5&7200&$(\A_5\times\A_5):\C_2$\\
&6&8160&$L_{16}:\C_2$\\
$[5;3,3]$&14&3840&$(E_5:\A_5)\times\C_2$\\
&19&3840&$E_6:\A_5$\\
&20&7680&$(E_6:\A_5)\times\C_2$\\\bottomrule
\end{tabular}
\caption{Actual tuples in the immutable eight-type core.}\label{tab}
\end{table}
For the three $E_d:\A_5:\C_2$ square rows the checker in fact constructs
an $\Sym_5$ complement. Its normal alternating subgroup and a transposition
give the displayed successive split extensions. The two complements of
order36 are explicitly $\Sym_3\times\Sym_3$, so their displayed two-step
series also follows, with the actual induced actions.

For $X_{32}$ the checker finds $|X_{32}|=32$ and
$Z(X_{32})=X_{32}'$ of order2, with generator squares in this central
subgroup. The quotient is elementary abelian: its generators commute and
have order at most2. Thus $X_{32}^2X_{32}'=Z(X_{32})$; the elementary
Frattini formula for a finite2-group shows that the Frattini subgroup also
equals this order2 group. This proves the asserted extraspecial property;
no sign or small-group identifier absent from the source is inferred.

For the affine rows we construct translations of $\mathbb F_2^3$ and the
literal natural $\mathrm{GL}_3(2)$ action on its eight vectors. A retained
bijective generator homomorphism identifies the linear group with $L_7$.
The resulting affine group has order1344. The required direct products
are then named models, compared by explicit bijective homomorphisms to
the entire actual $H$. No group is identified merely by its order.
\newpage
\section{The central nonsplit row and execution interface}
For $(\tau,j)=((3,3,3,4),12)$ the certificate gives
\[
 D\triangleleft H,\quad |D|=3840,\quad [H:D]=2,\quad
 Z=\C_4\leq Z(D),\quad Z\cap D'\ne1.
\]
It also gives $N\triangleleft D$ of order64 containing $Z$, and a subgroup
$U\leq D$ with a surjective homomorphism $U\to\A_5$ whose kernel is exactly
$Z$. Generator squares and commutators of $N$ lie in $Z$, so
$N/Z\cong E_4$. The checked equalities $N\cap U=Z$ and $NU=D$ give
$D/Z\cong E_4:\A_5$. If the central extension by $Z$ split, then
$D=Z\times K$ for a complement $K$, and $D'=K'$ would be disjoint from
$Z$. This contradicts $Z\cap D'\ne1$, proving the dot assertion.
An actual involution outside $D$ generates $H$ with $D$, proving the final
split factor. The other split $L_{16}$ rows additionally satisfy
$C_H(D)=1$ for their normal $D\cong L_{16}$; their external action is
therefore faithful, rather than an unproved identification by order.

\paragraph{Finite verification.}
The three certificate files contain only literal permutations, specified
subgroup generators, homomorphism images and outside elements. The verifier
reads the original core tuple and reconstructs every group from generators.
It checks containment, exact orders, normality, intersections, elementary
abelian tests, each claimed full generated group, and actual bijectivity
of each generator homomorphism. The twelve series records, nine direct-model
records and one central-extension record have exactly the22 distinct keys
in Table~\ref{tab}. Five damaged certificates independently corrupt a model
image, an elementary kernel, an outside involution, the order4 central
kernel, or the linear-group bridge; each is rejected.

These are direct positive certificates. Their earlier search routines are
candidate providers only, and their execution is not a completeness premise.
The accepted core separately supplies the necessary incidence mathematics
and complete normal-quotient searches. Default replay checks all22 groups
and the five damaged inputs; it does not rerun those searches. Python
optimization is rejected before output. Complete native sources are included.

\paragraph{Sources, software and scope.}
The source is M.~Fernandes, D.~Leemans and A.~I.~Weiss, the fixed sixteen-page
L033 article, Table~3 on physical page14. Its positive group entries are the
objects recovered here; we make no priority claim. The exact permutation,
homomorphism and finite-field calculations use GAP4.15.1. Its documented
$\mathrm{PSL}(2,q)$ constructor is the action of determinant-one matrices
on projective lines, and the standard alternating/symmetric constructors
are named-model interfaces. These software interfaces are retained explicitly,
not formally verified by this computation. The affine model itself is built
from the stated vector action. The acquired LINS completeness and finite
classification premises are unchanged and are not used for a new census here.

The project root supplied the original normal-quotient probe; the publication
agent authored these group certificates, verifier and companion proof.
Earlier distinct core reviews are acquired with their original roles.
The two source orders336 and2160 are excluded by that core, not by this
naming argument. The ninth diagram, its two printed orders3840 and9720,
and all remaining whole-source obligations are outside this theorem.
\end{document}
